% TOP_PROBLEM: {"id": 1100203, "title": "Questions in Geometric Group Theory — Q 2.3", "category_id": 17, "category": {"id": 17, "name": "group_theory", "display_name": "Group Theory", "description": "Problems about groups, group actions, representations, and related algebraic structures.", "slug": "group-theory", "order_index": 17, "created_at": "2026-07-31T15:26:25.671Z"}, "status": "open", "problem_number": "AMR-010-0203", "set_id": 13}
% FIRST_POSTED: 2026-10-06
% ATTEMPT_STATUS: solved
% OPENAI_CLAIM: {"schema_version":1,"source_repo":"https://github.com/openai/math","source_commit":"adc7f1241b42e322a6451854ab7e4b4c146bf78a","release_date":"2026-10-06","match":"full","resolution":"proved","independently_reviewed":false,"second_pass":"agreed","families":[{"family":"249","title":"A finitely generated Eilenberg–Ganea counterexample","subject":"Group theory","manuscripts":[{"title":"A finitely generated counterexample to the Eilenberg–Ganea conjecture","pdf_path":"preprints/A-finitely-generated-counterexample-to-the-Eilenberg-Ganea-conjecture-September-23-2026/paper.pdf","readme_path":"preprints/A-finitely-generated-counterexample-to-the-Eilenberg-Ganea-conjecture-September-23-2026/README.md","date":"2026-09-23","theorem_ref":"thm:main (sections/introduction.tex:15)","statement_sources":[{"path":"preprints/A-finitely-generated-counterexample-to-the-Eilenberg-Ganea-conjecture-September-23-2026/build/sections/introduction.tex","line":15,"label":"thm:main"}]}],"lean":{"doc_path":"lean/docs/249.md","comparators":["lean/ComparatorChallenges/EilenbergGanea.json","lean/ComparatorChallenges/EilenbergGanea.lean"],"declarations":[],"covers_main_theorem":true},"reasoning_trace":null}]}
\documentclass[11pt]{article}
\usepackage[margin=25mm]{geometry}
\usepackage{fontspec}
\setmainfont{Latin Modern Roman}
\usepackage{amsmath,amssymb,mathtools}
\usepackage{unicode-math}
\setmathfont{Latin Modern Math}
\usepackage{enumitem}
\usepackage{xurl,hyperref,bookmark}
\hypersetup{colorlinks=true,urlcolor=blue}
\setcounter{secnumdepth}{0}
\setlength{\parindent}{0pt}
\setlength{\parskip}{5pt}
\setlength{\emergencystretch}{3em}
\begin{document}
\section*{Questions in Geometric Group Theory — Q 2.3}

\subsection{OpenAI's claim}
OpenAI presents a proof in the selected manuscript. A finitely generated residually finite group is claimed to have integral cohomological dimension two but no two-dimensional classifying space.

\subsection{Mathematical statement of the claimed result}
\paragraph{Result 1.} Write \(\operatorname{cd}\) and \(\operatorname{gd}\) for integral cohomological dimension and the minimum dimension of a classifying space. There is a group \(G\), residually finite and finitely generated, for which
\[
 \operatorname{cd}G=2,\qquad\operatorname{gd}G=3.
\]
It takes a height kernel in a right-angled Artin group built from an acyclic flag triangulation of the presentation complex
\[
 \langle x,y\mid x^2=y^5,\ x^2=(xy^{-1})^3\rangle.
\]
Its nonexistence conclusion applies to every two-dimensional \(K(G,1)\), including those with infinitely many cells.

\subsection{Outline of the argument}
Degree estimates on conjugation-invariant parts of word fibers obstruct arbitrary aspherical presentations. Exact boundary-chain counts and quadratic bounds compare any candidate low-dimensional model with the constructed group, forcing geometric dimension three despite the shorter cohomological resolution.

\subsection{Scope relative to this problem}
The main theorem provides exactly the requested existence example; finite generation and residual finiteness are additional properties. Within these assumptions, this is a claimed affirmative answer to the stated question.

\subsection{Formal verification}
For family 249, the documented example is finitely generated and residually finite, has integral cohomological dimension two, and has geometric dimension three. The listed Comparator files are \texttt{EilenbergGanea.lean}. The inventory does not supply a named declaration for these Comparator entries. This site has not executed the Lean code. The scope descriptions do not certify the remaining manuscript arguments.

\subsection{Status caveat}
OpenAI attributes these manuscripts to work by an internal model. Its release cautions that arguments without formalization can contain mistakes. This catalogue has not independently checked the proofs; the labels report the mathematical claims attributed to OpenAI.

\subsection{Sources}
\begin{itemize}
\item OpenAI, \emph{A finitely generated counterexample to the Eilenberg–Ganea conjecture}, 2026-09-23; citation key \texttt{OAI:A-finitely-generated-counterexample-to-the-Eilenberg-Ganea-conjecture-September-23-2026}. \href{https://github.com/openai/math/blob/adc7f1241b42e322a6451854ab7e4b4c146bf78a/preprints/A-finitely-generated-counterexample-to-the-Eilenberg-Ganea-conjecture-September-23-2026/paper.pdf}{PDF (GitHub)}; \href{https://github.com/openai/math/raw/adc7f1241b42e322a6451854ab7e4b4c146bf78a/preprints/A-finitely-generated-counterexample-to-the-Eilenberg-Ganea-conjecture-September-23-2026/paper.pdf}{PDF (direct)}; \href{https://github.com/openai/math/blob/adc7f1241b42e322a6451854ab7e4b4c146bf78a/preprints/A-finitely-generated-counterexample-to-the-Eilenberg-Ganea-conjecture-September-23-2026/README.md}{Manuscript page and citation}; \href{https://github.com/openai/math/blob/main/preprints/A-finitely-generated-counterexample-to-the-Eilenberg-Ganea-conjecture-September-23-2026/paper.pdf}{Latest version}; Statement source: \href{https://github.com/openai/math/blob/adc7f1241b42e322a6451854ab7e4b4c146bf78a/preprints/A-finitely-generated-counterexample-to-the-Eilenberg-Ganea-conjecture-September-23-2026/build/sections/introduction.tex\#L15}{thm:main}.
\item Family 249: \href{https://github.com/openai/math/blob/adc7f1241b42e322a6451854ab7e4b4c146bf78a/lean/docs/249.md}{Lean scope}.
\item \href{https://github.com/openai/math/blob/adc7f1241b42e322a6451854ab7e4b4c146bf78a/lean/ComparatorChallenges/EilenbergGanea.json}{Comparator EilenbergGanea.json}.
\item \href{https://github.com/openai/math/blob/adc7f1241b42e322a6451854ab7e4b4c146bf78a/lean/ComparatorChallenges/EilenbergGanea.lean}{Comparator EilenbergGanea.lean}.
\item \href{https://github.com/openai/math/blob/adc7f1241b42e322a6451854ab7e4b4c146bf78a/CONTENTS.md}{OpenAI catalogue entry}.
\end{itemize}
\end{document}
